% Cover letter using letter.cls
\documentclass{letter}
\topmargin=-1in
\textheight=8in
\oddsidemargin=0pt
\textwidth=6.5in

\begin{document}

\signature{N. Ohayon Rozanes}
\longindentation=0pt
\let\raggedleft\raggedright

\begin{letter}{}

  \begin{center}
    {\large\bf N. Ohayon Rozanes} \\
    {611 Granbury Dr. \\ Allen, Texas, 75013 \\ 469-560-0141}
  \end{center} \vfill

  \opening{Dear Hiring Manager:}

  \noindent I am writing to apply for the Quantitative Technologist C++
  Intern role at Radix. I am pursuing a double degree in Data Science
  and Mathematics at the University of Texas at Dallas, and I have spent
  most of my time building performance-critical systems in C and C++. I
  am drawn to Radix because it is one of the few places where mathematical
  thinking and systems engineering are treated as equally serious disciplines
  rather than separate concerns handed off between teams.

  \noindent The bulk of my technical experience has been writing software
  that has real, immediate consequences if it misbehaves. At Exolambda I
  wrote real-time signal processing and control code for avionics hardware,
  where latency targets were non-negotiable and failures were visible on the
  bench immediately. That kind of feedback loop suits me well. At UTD I
  developed high-performance C++ search algorithms for a combinatorics
  research problem, which required genuinely understanding the language at
  the level of cache behavior, pruning strategy, and algorithmic complexity
  to make experiments tractable at scale. In both cases the work was
  fundamentally about squeezing correct behavior out of constrained systems,
  which is also how I understand what a trading system demands.

  \noindent I built a quantitative trading project, TradeX, that found
  correlated equity pairs through network analysis and backtested ML-driven
  strategies in simulation. It outperformed the naive baseline, but more
  importantly it made me want to work on something much closer to the real
  thing. I am not particularly interested in academic exercises for their
  own sake. Radix's description of taking a hacker's approach to testing
  ideas, dropping projects that aren't paying off, and pushing production
  code daily is exactly the kind of environment I am looking for. The
  mathematical degree is not background decoration for me; I actively use
  probabilistic and statistical reasoning when thinking about any system,
  and working in a firm where that is the expected baseline rather than an
  occasional nice-to-have is genuinely appealing.

  \noindent I would welcome the opportunity to work on Radix's trading
  infrastructure and to be held to the standards of a team that takes
  both the engineering and the math seriously.

  \closing{Sincerely yours,}

\end{letter}

\end{document}
